\section{数列の極限}\label{節:数列極限}
前章では,長方形の幅を小さくするときに $n$ を大きくし,割線を接線へ近づけるときに $h$ を $0$ に近づけた.どちらも,変化する数がある値に近づく様子を調べている.しかし「近づく」とは正確には何を意味するのだろうか.まず,番号 $n=1,2,3,\ldots$ の順に並んだ数列 $a_1,a_2,a_3,\ldots$ から考えよう.

\begin{definition}{数列の収束}
$n$ を十分大きくすれば,\,$a_n$ と一定の実数 $\alpha$ との差を,どのように小さく指定した正の数よりも小さくでき,それ以降も小さいままであるとき,数列 $\{a_n\}$ は $\alpha$ に\textbf{収束する}という.
\[
 \lim_{n\to\infty}a_n=\alpha
 \qquad\text{または}\qquad a_n\to\alpha\quad(n\to\infty)
\]
と書き,\,$\alpha$ を極限値という.
\end{definition}

「それ以降も」という部分は大切である.たとえば $a_n=(-1)^n$ は $1$ を何度も取るが,その間に $-1$ も取るので,\,$1$ に収束しない.極限は,一部の項だけが近いという意味ではない.また,\,$a_n=\frac1n$ はどの項も $0$ ではないが,\,$0$ に収束する.極限値を実際に取る必要もないのだ.

定義を数式で書くと,「任意の $\varepsilon>0$ に対して,ある自然数 $N$ が存在し,\,$n\ge N$ なら $|a_n-\alpha|<\varepsilon$ となる」ということである.たとえば $a_n=\frac1n$ なら,\,$N>\frac1\varepsilon$ となる自然数を選べばよい.すると $n\ge N$ に対して
\[
 \left|\frac1n-0\right|=\frac1n\le\frac1N<\varepsilon
\]
が成り立つ.極限の記号は,このような誤差の条件を短く表している.

数列が限りなく大きくなる場合には $a_n\to+\infty$,限りなく小さくなる場合には $a_n\to-\infty$ と書く.ただし,\,$\infty$ は実数の極限値ではない.このような発散と,\,$(-1)^n$ のような振動による発散は区別する.

収束する数列については,和,差,積,商の極限をそれぞれの極限から計算できる.ただし,商を考えるときは,分母の数列の極限が $0$ でないとし,分母が $0$ でない項だけを考える.\footnote{分母の極限 $\beta$ が $0$ でなければ,十分大きな $n$ で分母は $0$ でなくなる.}この断りのもとで,\,$a_n\to\alpha$,\,$b_n\to\beta$ のとき
\[
 a_n+b_n\to\alpha+\beta,\qquad
 a_nb_n\to\alpha\beta,\qquad
 \frac{a_n}{b_n}\to\frac\alpha\beta
\]
となる.これらの極限法則と,次の挟み撃ちの原理は,極限の定義から証明できる基本性質として用いる.

\begin{theorem}{挟み撃ちの原理}
十分大きな $n$ について $a_n\le b_n\le c_n$ が成り立ち,\,$a_n\to\alpha$,\,$c_n\to\alpha$ なら,\,$b_n\to\alpha$ である.
\end{theorem}

実際,\,$a_n,c_n$ がともに $\alpha$ との差を $\varepsilon$ より小さくすれば,その間の $b_n$ も $\alpha-\varepsilon<b_n<\alpha+\varepsilon$ に入る.前章の長方形による面積の計算は,\,$L_n\le S\le U_n$ の真ん中が $n$ によらない場合に,この原理を使う議論だった.

\begin{example}{極限の計算}
次の極限を求めよ.
\[
 \text{(1)}\quad\lim_{n\to\infty}\frac{3n^2-n+2}{2n^2+5},
 \qquad
 \text{(2)}\quad\lim_{n\to\infty}\bigl(\sqrt{n^2+3n}-n\bigr).
\]
\end{example}
\paragraphbf{方針}
\begin{enumerate}[label=(\arabic*)]
 \item 分子と分母を最も高い次数 $n^2$ で割り,分母の極限が $0$ でない形にして商の極限法則を使う.
 \item 共役な式 $\sqrt{n^2+3n}+n$ を分子と分母に掛けて,大きい数同士の差を商に書き換える.
\end{enumerate}
\paragraphbf{解答}
\begin{enumerate}[label=(\arabic*)]
 \item $\frac32$
 \item $\frac32$
\end{enumerate}
\paragraphbf{解法}
\begin{enumerate}[label=(\arabic*)]
 \item 分子も分母も大きくなるため,そのままでは商の極限法則を使えない.分子と分母を $n^2$ で割ると,
 \begin{align*}
  \frac{3n^2-n+2}{2n^2+5}
  &=\cfrac{3-\cfrac1n+\cfrac2{n^2}}{2+\cfrac5{n^2}}\\
  &\longrightarrow\frac32
 \end{align*}
 となる.分母の極限が $2\ne0$ となったので,商の法則を使えたのである.
 \item $\sqrt{n^2+3n}$ と $n$ がともに大きくなることだけでは,差の極限は決まらない.共役な式を分子と分母に掛けて差を商に変えると,
 \begin{align*}
  \sqrt{n^2+3n}-n
  &=\frac{(n^2+3n)-n^2}{\sqrt{n^2+3n}+n}\\
  &=\frac{3n}{\sqrt{n^2+3n}+n}\\
  &=\cfrac3{\sqrt{1+\cfrac3n}+1}\\
  &\longrightarrow\frac32
 \end{align*}
 となる.最後の行では平方根の連続性も使った.連続性の意味と基本関数の性質は\sectionref{節:関数極限連続}で説明する.
\end{enumerate}
(2) では,大きい数同士の差が有限の値へ近づいた.実際,\,$n=10,100,1000$ で $\sqrt{n^2+3n}-n$ を計算すると,およそ $1.402$,\,$1.489$,\,$1.499$ となり,\,$\frac32$ に近づいていく様子が見える.

\subsection{単調で有界な数列}
数列が\textbf{上に有界}であるとは,すべての $n$ について $a_n\le M$ を満たす一定の実数 $M$ が存在することをいう.下に有界であることも,同様に一定の下限によって定義する.増加を続けていても,上に有界なら無限に大きくなることはできない.

\begin{theorem}{単調有界数列の収束：採用する事実}
単調増加する数列が上に有界なら,有限の極限値へ収束する.単調減少する数列が下に有界なら,同じく有限の極限値へ収束する.
\end{theorem}

この定理の完全な証明には,実数の連続性に関する性質が必要になる.本書では実数の基礎理論までさかのぼらず,この定理を採用する.ただし,適用するときに単調性と有界性の両方を確かめることは省略しない.収束すると予想した式に極限値を代入するだけでは,収束の証明にはならないからである.



%
